%% =====================================================================
%%  B4-T-19-01 -- what B4 contributes to "The Uninvited Debt"
%%  -------------------------------------------------------------------
%%  LAYER A: APPEND-ONLY. Written once, then left alone. Material found
%%  only here sits in the `unique` slot and is protected by rule R10.
%% =====================================================================
%% @kind: source
%% @id: B4-T-19-01
%% @book: B4
%% @topic: T-19-01
%% @locator: ch. 2, pp. 43--86
%% @status: distilled
%% @unique: yes
%% @integrated: yes
%% @updated: 2026-09-19

\dossier{B4}{T-19-01}{\bookshort{B4}}{ch. 2, pp. 43--86}

\begin{distilled}
  \item The obligation to reciprocate is triggered by the receipt of a benefit, not by agreement to receive it; uninvited favours create a genuine sense of debt.
  \item The rule is enforced by the social cost of being seen not to reciprocate, which is why opting out is difficult and why the discomfort is about self-image rather than arithmetic.
  \item The debt can be drawn on for a request wholly unrelated to the original benefit, which separates it from exchange.
  \item Field settings in which a benefit is delivered so that public refusal is more costly than acceptance.
\end{distilled}

\begin{unique}
  \item The absence of consent as the load-bearing feature: the obligation arises from occurrence, not from agreement. Neither B1 nor B3 isolates this.
  \item The social-marking account of why the rule is hard to refuse, as distinct from a claim about gratitude as sentiment.
\end{unique}
